% Generated draft: requires mathematical and source review.
\documentclass[11pt]{article}
\usepackage[margin=25mm]{geometry}
\usepackage{fontspec}
\setmainfont{Latin Modern Roman}
\usepackage{amsmath,amssymb,mathtools,unicode-math}
\setmathfont{Latin Modern Math}
\usepackage{enumitem,needspace,xurl,hyperref}
\setlength{\parindent}{0pt}
\setlength{\parskip}{5pt}
\setlength{\emergencystretch}{3em}
\newcommand{\N}{\mathbb{N}}
\newcommand{\Z}{\mathbb{Z}}
\newcommand{\Q}{\mathbb{Q}}
\newcommand{\R}{\mathbb{R}}
\newcommand{\C}{\mathbb{C}}
\newcommand{\F}{\mathbb{F}}
\newcommand{\A}{\mathbb{A}}
\newcommand{\PP}{\mathbb{P}}
\newcommand{\E}{\mathbb{E}}
\newcommand{\eps}{\varepsilon}
\newcommand{\dd}{\,\mathrm{d}}
\begin{document}
\section*{UnsolvedMath ID: 1200015}
\textbf{Source problem code:} catalogue entry
\subsection{Definitions and mathematical statement}
Use finite nonempty simple graphs $G_n$. Write $N_n=|V(G_n)|$ and let $P_{G_n}(q)$ be the chromatic polynomial, whose value at a positive integer $q$ counts proper $q$-colorings. Listing its $N_n$ complex roots with algebraic multiplicity as $\lambda_{n,1},\ldots,\lambda_{n,N_n}$, define the chromatic-root probability measure
\[
\mu_n=\frac1{N_n}\sum_{j=1}^{N_n}\delta_{\lambda_{n,j}}.
\]
For the local-convergence condition, choose a uniformly random vertex $v_n$ of $G_n$. For each radius $r\geq0$ and each finite rooted graph $(H,o)$, require that
\[
\mathbb P\big((B_r(G_n,v_n),v_n)\cong(H,o)\big)
\]
converge as $n\to\infty$, with the limiting probabilities forming a probability distribution on rooted radius-$r$ balls. Here $B_r$ is the induced graph on vertices at graph distance at most $r$. This is the locally finite local-weak-convergence reading. The source does not say whether a uniform bound on all vertex degrees is part of the term locally convergent; that extra hypothesis must be specified rather than silently inserted.

\textbf{Question.} For every locally convergent sequence $(G_n)$ in the intended class and every nonnegative integer $k$, does the finite complex limit
\[
\lim_{n\to\infty}\int_{\mathbb C}z^k\,d\mu_n(z)
=\lim_{n\to\infty}\frac1{N_n}\sum_{j=1}^{N_n}\lambda_{n,j}^{\,k}
\]
exist? These are holomorphic moments, not absolute moments or mixed moments $z^k\overline z^{\ell}$. The supplied research summary reports partial work on chromatic measures but does not verify the fully general assertion as settled.
\subsection{Short English statement}
Does convergence of local graph neighborhoods force convergence of every normalized power sum of chromatic roots? The intended degree-bound convention is essential to identifying the exact statement.
\subsection{Sources}
\begin{itemize}
\item M. Ab\'ert, \emph{Some questions} (2010), Question 15, \url{https://www.renyi.hu/~abert/questions.pdf}; manual transcription, accessed 2026-07-29. The record says the 2010 author list was inspected, with current status not independently verified.
\item The supplied research summary names Csikv\'ari and collaborators for work on chromatic measures and their limits, but supplies no title, date, theorem, or bibliographic identifier.
\item Local registry \texttt{data/unsolvedmath/problems.json}, permanent ID $1200015$, source code catalogue entry, \emph{Some Questions --- Question 15}. The record labels this partially solved and PARTIAL-PROGRESS, with source-status and rights qualifications marked NEEDS\_REVIEW.
\end{itemize}
\end{document}
